\section{Normal Duality and Geometrization}
\label{sec:normal_duality}

\paragraph{From fixed-\(\beta\) reflection to geometric reflection.}
The map \(g\mapsto g^\vee\) at fixed \(\beta\) already reflects the scalar coordinate
\(y=\log\chi\).
A natural extension is to ask whether both \(\beta\) and \(g\) can change simultaneously so that
partner points are connected by a direction normal to the crossover curve.

\paragraph{Crossover manifold as a geometric object.}
At fixed \(\theta\), define
\begin{equation}
  y(\beta,g,\theta)=\log\chi(\beta,g,\theta),
  \qquad
  \mathcal{C}_\theta:=\{(\beta,g):y=0\}.
\end{equation}
In one-channel scaling,
\begin{equation}
  y=2\log g+\log\chi_0(\beta,\theta),
  \qquad
  \nabla y=\left(\partial_\beta\log\chi_0,\frac{2}{g}\right).
\end{equation}
So \(\nabla y\) defines the local normal direction to the self-dual contour.

\paragraph{Normal dual representative.}
Given a point \(p=(\beta,g)\) near \(\mathcal{C}_\theta\), project it to the closest crossover
point \(\pi(p)\) and follow the local normal ray through \(\pi(p)\).
The normal-dual partner \(p^\perp=(\beta^\perp,g^\perp)\) is defined as the point on the same ray
satisfying
\begin{equation}
  y(p^\perp)=-y(p).
\end{equation}
Thus normal duality is the geometric analogue of strong-weak inversion, now in two coordinates.

\paragraph{What this adds physically.}
The fixed-\(\beta\) involution is recovered when the normal is nearly vertical.
Away from that limit, normal duality encodes how temperature and coupling should be co-varied to keep
branch distance symmetric around crossover.
This is useful when experiments or numerics sweep both controls simultaneously.

\paragraph{Practical construction recipe.}
For numerical implementation:
\begin{enumerate}
  \item compute the crossover contour \(\mathcal{C}_\theta\) from \(\chi=1\),
  \item estimate local normal \(\hat\nu\propto\nabla y\),
  \item for each point \(p\), solve along the normal ray for \(p^\perp\) with \(y(p^\perp)=-y(p)\).
\end{enumerate}
This produces a simultaneous \((\beta,g)\)-dual map without introducing any additional model
assumption beyond smoothness of \(y\).

\paragraph{Formal statement.}
A local existence/involution theorem for this map is stated and proved in
the Appendix.
